\section{General number-field annular transport on the original flat quotient}
\label{sec:general-annular-transport}

\paragraph{Status.}
The complete author reading and two independent
whole-proof readings have passed.  It addresses the pre-existing question
\texttt{Q-MA-LA-GENERALS}.  The field, source topology, Fourier
operator and two physical norm cutoffs below are part of the statement.
They are not replaced by a radial additive model.

Let $K$ be a fixed number field of signature $(r_1,r_2)$ and degree
$n=r_1+2r_2$.  Let $S$ be a finite structural set of places containing
all infinite places.  Here structural means Meyer's sufficiently
large condition
$C_K/\mathcal O_{\overline S}^{\times}\simeq C_S$; no prime in $S$
is required to be principal.  Set
\[
 G=\mathbb A_S^\times,\qquad
 \Gamma=\mathcal O_{K,S}^\times,\qquad
 G^1=\ker|\cdot|_S,\qquad C_S=G/\Gamma,
 \qquad C_S^1=G^1/\Gamma.
 \tag{MA-GA.1}
\]
The $S$-unit group is discrete and cocompact in $G^1$.
In particular $C_S^1$ is compact, including when the ordinary
unit rank is positive.  Use the positive homomorphic norm splitting
$\nu(e^x)$ whose infinite coordinates are $e^{x/n}$ and whose
finite coordinates are $1$.  The normalized absolute value at a
complex place is the square of the ordinary absolute value, so
$|\nu(e^x)|_S=e^x$.  Give $C_S^1$ probability Haar measure.

The original flat source, with its LF topology, is
\[
 V_S=\mathcal S_{00}(K_\infty)
       \mathbin{\widehat\otimes}_{\iota}
       \bigotimes_{v\in S_f}C_c^\infty(K_v^\times).
 \tag{MA-GA.2}
\]
Here every ordinary partial derivative of the infinite component
vanishes on each coordinate hyperplane $x_v=0$.  At a finite
place the factor means a locally constant, compactly supported
function whose support avoids zero, extended by zero to $K_v$.
It does not mean the larger multiplicative Schwartz space with
unbounded valuation support.  A finite-place LF stage fixes a
compact punctured support and an open invariance subgroup; its
finite-place factor is finite dimensional.  We specify
the topology by these strict LF stages, as in MA-LA.41;
no interchange of a projective tensor product with arbitrary
inductive limits is assumed.
Intrinsically this is the coordinate-flat closed subspace
of the additive Schwartz--Bruhat space
$\mathcal S(\mathbb A_S)$ with its inherited LF topology.
The inductive tensor notation in (MA-GA.2) records that
topology.  We do not assert that the printed projective
tensor symbol in Masullo's (5.3.4.1), taken as a separate
topological prescription, equals it.

Fix a smooth unitary character $\chi$ of $C_S^1$ and extend it
trivially along $\nu$.  Its raw compact type is $\chi^{-1}$.
In the centered norm coordinate the corresponding Hilbert space
is $H_\chi=L^2(\mathbb R,dx)$.  A scalar function $F$ represents
$\chi(\bar k)^{-1}F(x)$, where $k\in G^1$ and $x=\log|t|_S$.
Write
\[
 \mathcal E=\{F\in C^\infty(\mathbb R):
      \sup_x e^{N|x|}|F^{(j)}(x)|<\infty
          \text{ for every }N,j\geq0\}.
 \tag{MA-GA.3}
\]
All character and norm identifications below are made on this
actual quotient, with the above Haar normalization.

Choose $A,B\in C^\infty([0,\infty))$, between zero and one,
equal to one on $[0,1]$ and zero on $[2,\infty)$.  For
$X\geq2\epsilon>0$, $Y>0$, and $g\in C_c^\infty(\mathbb R)$,
let
\[
\begin{aligned}
 a_{\epsilon,X}(x)&=A(e^x/X)-A(e^x/\epsilon),\\
 (C_gF)(x)&=\int_{\mathbb R}g(y)F(x-y)\,dy,\\
 Q_{\chi,Y}&=\mathfrak S_{\bar\chi}
               M_{B(e^x/Y)}\mathfrak S_\chi,\\
 T^\chi_{X,Y,\epsilon}(g)&=M_{a_{\epsilon,X}}C_gQ_{\chi,Y}.
\end{aligned}
 \tag{MA-GA.4}
\]
The operators $\mathfrak S_\chi:H_\chi\to H_{\bar\chi}$ are
the actual descended additive Fourier transforms with the
coherent local additive characters and self-dual measures of
MA-NR-001/MA-CQ.  Thus $B(e^x/Y)$ cuts the physical total norm
on the Fourier side.  Neither this cutoff nor $A(e^x/X)$ is
a cutoff in an individual archimedean coordinate.

Let $RF(x)=F(-x)$, interchanging the two compact sectors, and put
\[
 U_\chi=\mathfrak S_{\bar\chi}R,
 \qquad \mathcal M(U_\chi F)(z)=m_\chi(z)\mathcal MF(z),
 \qquad \mathcal MF(z)=\int_{\mathbb R}F(x)e^{zx}\,dx.
 \tag{MA-GA.5}
\]
This defines $m_\chi$ by the actual local Fourier operator;
all its conductor, different and phase factors are retained.

\paragraph{MA-GA-001 (general structural-$S$ source, lift and annular rate).}
The separated $\Gamma$-coinvariants of (MA-GA.2) identify
topologically with the smooth functions on $C_S$ having every
exponential norm weight.  Their fixed $\chi^{-1}$ compact
sector is precisely (MA-GA.3).  Centered periodization has a
continuous section into the original source; its central
Hilbert completion is the stated $H_\chi$.

The operator $T^\chi_{X,Y,\epsilon}(g)$ in (MA-GA.4) is trace class on that completion
and has a continuous nuclear lift into (MA-GA.2).  The lift
and its quotient operator have the same trace, computed from
one nuclear series.  The intermediate Fourier vectors are
allowed in the actual completion and need not be flat.

For $g(0)=0$, define the actual local trace
\[
 \mathcal W_{S,\chi}(g)
    =-\frac1{2\pi}\int_{\mathbb R}
             \mathcal Mg(it)\frac{m_\chi'(it)}{m_\chi(it)}\,dt.
 \tag{MA-GA.6}
\]
For each $0<\delta<1/2$ there is a constant $C_\delta$,
depending on the fixed field, $S,\chi,A,B,g,\delta$ but not
on $X,Y,\epsilon$, such that, with $U=XY$ and $V=\epsilon Y$,
\[
\begin{aligned}
 \left|\operatorname{Tr}T^\chi_{X,Y,\epsilon}(g)
                    -\mathcal W_{S,\chi}(g)\right|
   &\leq C_\delta(U^{-\delta}+V^\delta)
          &&(U\geq1,\ V\leq1),\\
 \left|\operatorname{Tr}T^\chi_{X,Y,\epsilon}(g)\right|
   &\leq C_\delta V^{-\delta}
          &&(V\geq1).
\end{aligned}
 \tag{MA-GA.7}
\]
For tests in a fixed compact support the constants are
controlled by finitely many smooth seminorms of $g$.
Consequently the two coupled limits are, respectively,
$\mathcal W_{S,\chi}(g)$ when $XY\to\infty$ and
$\epsilon Y\to0$, and zero when $\epsilon Y\to\infty$.
These statements allow a nontrivial class group, every unit
rank and every archimedean signature.  They use no global
zero-location hypothesis.

\begin{proof}
\emph{1. The original flat topology and its periodization.}
For a real coordinate let $r_v$ be its ordinary absolute
value, and for a complex coordinate let it be its ordinary
modulus.  Restriction of the flat archimedean source to
$K_\infty^\times$ is characterized, for every $N\geq0$
and every invariant differential operator $D$, by
\[
 \sup_{t\in K_\infty^\times}
    \exp\!\left(N\sum_{v\mid\infty}|\log r_v(t)|\right)
       |D f(t)|<\infty.
 \tag{MA-GA.8}
\]
These seminorms and the original flat subspace topology are
equivalent.  Here is a direct proof of the estimates needed
in both directions.  Split the coordinates into $r_v\leq1$
and $r_v>1$.  In each small coordinate apply Taylor's formula
along the segment from zero to that coordinate, to any
prescribed order.  Every lower Taylor coefficient is zero
on the entire coordinate hyperplane, also after differentiating
in the other coordinates.  Iterating gives a product of
arbitrarily high powers of the small $r_v$.  A sufficiently
large ordinary Schwartz weight bounds the product of negative
powers of all large $r_v$.  Invariant derivatives are finite
sums of polynomial coordinate coefficients times ordinary
derivatives, so the same argument applies after such a
derivative.  Each seminorm in (MA-GA.8) is therefore bounded
by finitely many higher ordinary Schwartz seminorms.

Conversely ordinary derivatives on the punctured coordinates
are combinations of invariant derivatives with coefficients
bounded by powers of $r_v^{-1}$.  The estimates (MA-GA.8)
imply that every such derivative tends to zero faster than
every power at each coordinate hyperplane, uniformly when
the other coordinates range over a compact set.  Extending
by zero gives a smooth function there: in one real direction
use the fundamental theorem of calculus on the segment from
zero, then repeat for every derivative; in a complex coordinate
use the same argument in its two real directions.  Repeating
over the coordinates proves the asserted simultaneous flatness.
The large-coordinate estimates give all ordinary Schwartz
seminorms.  This proves the converse continuity as well.

It also proves density of $C_c^\infty(K_\infty^\times)$ in
the flat source.  Indeed multiply by smooth cutoffs that equal
one on $|\log r_v|\leq R$ and vanish on
$|\log r_v|\geq R+1$, for every $v$.  The error tends to zero
in each seminorm (MA-GA.8), using a strictly higher exponential
weight.  At finite places this procedure changes neither
the support nor the invariance level.  Thus
$C_c^\infty(G)$ is dense in $V_S$, with convergence in each
original finite-place LF stage.

Let $\mathcal E(C_S)$ mean the analogous all-exponential-weight
smooth space on $C_S\simeq C_S^1\times\mathbb R$, with the
usual LF union over finite compact invariance levels.
Define centered periodization by
\[
 (\mathcal P f)(\bar t)
      =|t|_S^{1/2}\sum_{\gamma\in\Gamma}f(\gamma t).
 \tag{MA-GA.9}
\]
The sum and all its differentiated, exponentially weighted
versions converge, and $\mathcal P:V_S\to\mathcal E(C_S)$
is continuous.  We verify the point that is absent from a
finite-unit-shell proof.  Cover the compact quotient $C_S^1$
by representatives in one compact subset of $G^1$.
On a fixed source LF stage and these representatives, the
finite valuations of a nonzero summand in (MA-GA.9) permit
only finitely many valuation vectors of $\gamma$.
Each nonempty such fiber is a coset of $\mathcal O_K^\times$.
The archimedean logarithms of ordinary units form a lattice
in the norm-zero real hyperplane, with finite kernel the
roots of unity.  This also follows by projecting their
discrete cocompact subgroup through the compact angular
kernel.  The lattice has polynomial counting growth.

For a representative $k$ in the fixed compact set and
$t=k\nu(e^x)$, the vector of archimedean logarithms of
$\gamma t$ is a fixed bounded vector plus its unit-lattice
vector plus $(x/n,\ldots,x/n)$.  Normalized complex logs
have the corresponding factors two.  Norm equivalence
between this direct decomposition and Euclidean coordinates
gives constants $c,C>0$ such that
\[
 \sum_{v\mid\infty}|\log r_v(\gamma k\nu(e^x))|
       \geq c\bigl(|x|+\|\ell(\gamma)\|\bigr)-C
 \tag{MA-GA.10}
\]
on each of the finitely many allowed valuation fibers.
The fiber representatives are absorbed in $C$.
An arbitrarily high weight in (MA-GA.8) now dominates the
factor $e^{x/2}$, every required norm weight and the lattice
sum.  Differentiation only replaces the source seminorm by
one of higher order.  This proves the assertion uniformly
on the compact quotient.  The finite-place invariance level
of the periodization is fixed by that of the source stage,
so this is continuity in the stated LF topology.

\emph{2. A section and the exact closed relation space.}
Because the action of the discrete subgroup $\Gamma$ on
$G^1$ is proper and its quotient is compact, choose
$b_0\in C_c^\infty(G^1)$, $b_0\geq0$, whose translates
cover the quotient with positive sum.  The sum is locally
finite.  The function
\[
 b(k)=\frac{b_0(k)}{\sum_{\gamma\in\Gamma}b_0(\gamma k)}
 \quad\text{satisfies}\quad
       \sum_{\gamma\in\Gamma}b(\gamma k)=1.
 \tag{MA-GA.11}
\]
It is smooth, compactly supported and has a fixed finite
invariance level.  The formula includes finite torsion;
no division by an unrecorded number of roots of unity occurs.
Define, for $\Psi\in\mathcal E(C_S)$,
\[
 (\sigma\Psi)(k\nu(e^x))=e^{-x/2}b(k)\Psi(\bar k,x).
 \tag{MA-GA.12}
\]
Then $\mathcal P\sigma=1$.  This section takes values in
the original $V_S$, continuously.  Indeed on its support
all archimedean moduli are comparable to $e^{x/n}$ and all
finite valuations are in a fixed finite set.  Every source
seminorm (MA-GA.8) is therefore controlled by finitely many
norm-exponential and compact-direction seminorms of $\Psi$.
The finite invariance level is the common refinement of
those of $b$ and $\Psi$.  The estimates just proved give
the flat zero extension, not merely a multiplicative
Schwartz function.  On compact norm support the image is
in $C_c^\infty(G)$ in a fixed original LF stage.

Let $L_\gamma f(t)=f(\gamma^{-1}t)$.  The kernel of
$\mathcal P$ is exactly
\[
 \ker\mathcal P
     =\overline{\operatorname{span}
          \{L_\gamma f-f:\gamma\in\Gamma,\ f\in V_S\}}^{\,V_S}.
 \tag{MA-GA.13}
\]
For a compactly supported $f$ put
$h_\gamma(t)=b(k)f(\gamma t)$, $t=k\nu(e^x)$.
Only finitely many $h_\gamma$ are nonzero: their possible
$t$ lie in a compact set, and properness of the action
then permits only finitely many $\gamma$.  Directly,
\[
 \sum_\gamma h_\gamma=\sigma\mathcal P f,
 \qquad
 \sum_\gamma L_\gamma h_\gamma=f,
 \qquad
 f-\sigma\mathcal P f
          =\sum_\gamma(L_\gamma-1)h_\gamma.
 \tag{MA-GA.14}
\]
Each term belongs to the original compact flat source.
For general $f$ use the compact approximations from step 1
and continuity of both maps.  If $\mathcal P f=0$ this
places $f$ in the closed relation space.  The converse
follows by reindexing the convergent periodization sum.
Equations (MA-GA.12--14) give the topological quotient
isomorphism, with continuous inverse induced by $\sigma$.

Compact-character projection on $C_S^1$, with its probability
Haar measure, is a continuous projector on this quotient.
Its $\chi^{-1}$ range is exactly $\mathcal E$ of (MA-GA.3).
In particular
\[
 (\sigma_\chi F)(k\nu(e^x))
       =e^{-x/2}b(k)\chi(\bar k)^{-1}F(x)
 \tag{MA-GA.15}
\]
is a continuous section of this sector.  Centering turns
$L^2(C_S,|t|_S\,d^\times t)$ on the sector into $L^2(dx)$.
Since $C_c^\infty(\mathbb R)\subset\mathcal E$, its completion
is exactly $H_\chi$.  This construction neither assumes a
trivial ideal class group nor chooses generators of primes.
Compact averaging occurs on $C_S^1$, not on the generally
noncompact group $G^1$.

The balanced splitting is a coordinate choice, not a change
of the physical operator.  If the original source uses
another positive homomorphic splitting, write
$\nu_0(e^x)=\kappa(x)\nu(e^x)$ and
$\chi(\kappa(x))=e^{i\lambda x}$.  The raw $\chi^{-1}$
convention gives $F_0(x)=e^{-i\lambda x}F(x)$.
With $D=M_{e^{-i\lambda x}}$ one has exactly
\[
 D C_gD^{-1}=C_{e^{-i\lambda\cdot}g},\qquad
 D M_{A(e^x/X)}D^{-1}=M_{A(e^x/X)},\qquad
 Q^{\nu_0}_{\chi,Y}=D Q^\nu_{\chi,Y}D^{-1}.
 \tag{MA-GA.15a}
\]
The Fourier-side identity follows by using the conjugate
coordinate change in the dual compact sector; the two
changes cancel around its physical norm multiplier.
Thus the theorem for every compact smooth $g$ applies
to the original splitting with the displayed test twist.
Its condition $g(0)=0$ is preserved.  Neither physical
norm cutoff has changed.

\emph{3. The actual local Fourier operator and its bounds.}
Use the additive Fourier transform on $\mathcal S(\mathbb A_S)$,
the chosen self-dual measures, and
$Jf(t)=|t|_S^{-1}f(t^{-1})$ on $G$.
Meyer's descent lemma applies to the equivariant operator
$\mathcal F J$: it is a continuous unitary multiplier on
$L^2(C_S,|t|_S\,d^\times t)$.  Thus the descended additive
Fourier operator interchanges the two compact sectors and
reverses the norm coordinate.  Since $-1\in\Gamma$, its
square on the quotient is one.  In the present notation,
\[
 \mathfrak S_{\bar\chi}\mathfrak S_\chi=1,
 \qquad U_\chi\text{ is unitary},\qquad
 \mathfrak S_{\bar\chi}=U_\chi R,
 \quad\mathfrak S_\chi=R U_\chi^{-1}.
 \tag{MA-GA.16}
\]
This is the descent of the actual transform on the source,
not a choice of an unrelated unitary with the same scalar
trace.  The flat quotient embeds into the full local
coinvariant space $\mathcal T(C_S)$ by (MA-GA.9).
Meyer's full local-coinvariant theorem identifies that space
with the separated coinvariants of $\mathcal S(\mathbb A_S)$.
Our additional argument (MA-GA.8--14) identifies the smaller
flat quotient with its own, generally stronger, topology.
No claim that Fourier preserves flatness is required.

For completeness the exact local functional equation can
define its multiplier without a Fourier-sign convention
ambiguity.  Define $\gamma_v(s,\omega_v,\psi_v)$ by
\[
 Z_v(\mathcal F_v f,\omega_v^{-1},1-s)
       =\gamma_v(s,\omega_v,\psi_v)Z_v(f,\omega_v,s).
\]
Then Mellin unfolding of (MA-GA.9), and $U_\chi=\mathcal FJ$
on its sector, give
\[
 m_\chi(z)=\prod_{v\in S}
    \gamma_v(1/2-z,\chi_v^{-1},\psi_v).
 \tag{MA-GA.17}
\]
The raw $\chi^{-1}$ convention is important here: its
coefficient is extracted by integrating against $\chi$, so
the unfolded Mellin transform is $Z(f,\chi,1/2+z)$.
For $s=1/2+z$, substitution $t\mapsto t^{-1}$ gives
$Z(Jf,\chi^{-1},1-s)=Z(f,\chi,s)$.  Applying the displayed
functional equation with the character $\chi^{-1}$ and
parameter $1-s$ therefore gives exactly (MA-GA.17), rather
than its reciprocal or its conjugate-character version.
Equivalently each factor is
\[
 \gamma_v(1/2-z,\chi_v^{-1},\psi_v)
   =\varepsilon_v(1/2-z,\chi_v^{-1},\psi_v)
       \frac{L_v(1/2+z,\chi_v)}
            {L_v(1/2-z,\chi_v^{-1})}.
 \tag{MA-GA.17a}
\]
The normalization of the quotient Haar measure cancels in
this ratio.  Here $\psi_v$ means the actual Fourier kernel,
including the conjugation specified in MA-CQ.

Explicitly at a real place write
$\chi_v(t)=\operatorname{sgn}(t)^{e_v}|t|_v^{i\tau_v}$.
At a complex place write
$\chi_v(t)=(t/|t|_{\mathrm{ordinary}})^{j_v}|t|_v^{i\tau_v}$.
Here $e_v\in\{0,1\}$, $j_v\in\mathbb Z$ and
$\tau_v\in\mathbb R$; $\chi_v^{-1}$ changes $\tau_v,j_v$
to their negatives and leaves $e_v$ unchanged.
At an unramified finite place let $\omega_v=\chi_v(\pi_v)$
for any local uniformizer.  Thus
$L_v(s,\chi_v)=(1-\omega_vq_v^{-s})^{-1}$ and
$\chi_v^{-1}(\pi_v)=\overline{\omega_v}$.
The real and complex local $L$ factors are respectively
$\Gamma_{\mathbb R}(s+i\tau_v+e_v)$ and
$\Gamma_{\mathbb C}(s+i\tau_v+|j_v|/2)$.

Tate's real and complex tables express the archimedean
factors, up to their actual exponential monomial and unit
phase, as
\[
 \frac{\Gamma_{\mathbb R}(1/2+z+i\tau_v+e_v)}
      {\Gamma_{\mathbb R}(1/2-z-i\tau_v+e_v)},
 \qquad
 \frac{\Gamma_{\mathbb C}(1/2+z+i\tau_v+|j_v|/2)}
      {\Gamma_{\mathbb C}(1/2-z-i\tau_v+|j_v|/2)}.
 \tag{MA-GA.18}
\]
Here $e_v\in\{0,1\}$, $j_v\in\mathbb Z$, and
$\tau_v\in\mathbb R$ are the fixed local parameters of
the unitary character.  The finite unramified factors have
the form
\[
 c_v e^{d_vz}
    \frac{1-\overline{\omega_v}q_v^{z-1/2}}
         {1-\omega_v q_v^{-z-1/2}},
 \qquad |c_v|=|\omega_v|=1,\quad d_v\in\mathbb R;
 \tag{MA-GA.19}
\]
the ramified factors are nonzero exponential monomials
with the actual Gauss phase.  Only their analytic bounds
are being read from these forms; (MA-GA.17) retains every
normalizing factor.  The numerator gamma poles lie on or
to the left of $\Re z=-1/2$; the finite denominator poles
lie on that line.  Denominator gamma poles give zeros,
not new poles.  Stirling's formula and its differentiated
forms therefore give
\[
 \begin{gathered}
 m_\chi\text{ holomorphic on }\Re z>-1/2,\qquad
 |m_\chi(it)|=1,\\
 m_\chi^{(j)}(c+it)\text{ has polynomial growth in }t
       \text{ on every fixed closed vertical substrip},\\
 (m_\chi^{-1})^{(j)}(c+it)\text{ has polynomial growth
       on every closed substrip of }-1/2<\Re z<1/2.
 \end{gathered}
 \tag{MA-GA.20}
\]
There is no uniformity claim as $K,S,\chi$ vary.  Polynomial
bounds for each derivative suffice below.

\emph{4. An exact kernel and its two weighted output estimates.}
Write $b_1(x)=B(e^x)$, $v(x)=1-b_1(-x)$ and
$T_v=M_v-U_\chi M_vU_\chi^{-1}$.  Thus $v=1$ far left,
$v=0$ far right, $v'$ is compactly supported, and
$\int v'=-1$.  The bounded identities (MA-GA.16) imply
\[
 Q_{\chi,1}=U_\chi M_{b_1(-x)}U_\chi^{-1}
                =M_{b_1(-x)}+T_v.
 \tag{MA-GA.21}
\]
Let $G(z)=\mathcal Mg(z)$ and
$P_v(z)=\int v'(x)e^{zx}\,dx$.  For $c>-1/2$ set
$H_c=L^2(\mathbb R,e^{2cx}dx)$ and use the unitary Fourier
map $F\mapsto(2\pi)^{-1/2}\mathcal M F(c+it)$ on $H_c$.
For $|c|<1/2$, the actual operator $K_g=C_gT_v$, viewed
from $H_0$ to $H_c$, has kernel
\[
 k_c(t,u)=\frac{G(z)}{2\pi}
       \frac{P_v(z-w)}{z-w}
       \left(\frac{m_\chi(z)}{m_\chi(w)}-1\right),
 \qquad z=c+it,\quad w=iu.
 \tag{MA-GA.22}
\]
We justify both the literal operator and the negative
output lines.  Start with input having smooth compact
Fourier support.  Its image under $U_\chi^{-1}$ is
Schwartz on the central line by (MA-GA.20).  If $c>0$,
multiplication by $v$ maps this vector into the Schwartz
space on line $c$.  The actual weighted Fourier extension
is compatible with the central one.  For compact inputs
this is the local functional equation on the compact
source section; for the present input, cut off in the
physical norm coordinate and pass to the limit in both
central and $c$-weighted Schwartz topologies.  The
polynomial multiplier bounds give convergence in the
weighted topology, and central unitarity identifies the
same local-distribution limit.  This is also the
compatibility in Meyer's local Fourier estimates.  Integration by parts gives
the convergent transform
\[
       \int_{\mathbb R}v(x)e^{(z-w)x}\,dx
                      =-\frac{P_v(z-w)}{z-w}
                     \qquad(c>0).
 \tag{MA-GA.23}
\]
Subtracting the two literal compositions in $T_v$ yields
(MA-GA.22), with the displayed sign.  There is no use of
the Fourier transform of a nonintegrable step as an
ordinary function on the central line.

The apparent singularity at $z=w$ in (MA-GA.22) is
removable: its quotient factor is
$m_\chi'(w)/m_\chi(w)$ there.  The bounds (MA-GA.20)
and the rapid vertical decay of $G(z)$ and $P_v(z-w)$
show that $k_c$ is Schwartz in $(t,u)$, uniformly for
$c$ in every compact subinterval of $(-1/2,1/2)$.
Near $z=w$ use the integral formula for the divided
difference; away from it use the polynomial bounds
directly.  The inequality
$1+|u|\leq(1+|t|)(1+|t-u|)$ shows that the two rapid
decays dominate every joint polynomial.  The same
argument applies after any $t,u$ derivative.

For the compact-Fourier-support inputs just used, the
right side of (MA-GA.22), integrated against the input,
is holomorphic in $z$ on this strip and rapidly decreasing
on every closed vertical substrip.  Shift the inverse
Mellin contour from a positive $c$ to any
$-1/2<c<1/2$.  The horizontal integrals tend to zero by
these bounds, and no pole is crossed.  Thus all these
kernels represent the same actual central output, not
different regularizations.  Density then identifies the
extended operator on all of $H_0$.

A Schwartz kernel on $\mathbb R^2$ is trace class.  One
direct proof expands it in the product Hermite basis:
integration by parts with both harmonic oscillators
gives coefficients decreasing faster than every power
in both indices.  Their absolute sum is finite, giving
a nuclear series of rank-one operators with summable
norms.  The same series permits the integral of the
diagonal, by the corresponding uniform Schwartz bounds.
Applied to (MA-GA.22), this proves
\[
 K_g:H_0\longrightarrow H_c\text{ is trace class}
                  \quad(-1/2<c<1/2).
 \tag{MA-GA.24}
\]
Its trace-class norm on any one line is bounded by
finitely many Schwartz kernel seminorms, hence by
finitely many seminorms of $g$ on a fixed compact
support.  On the central diagonal $P_v(0)=-1$, so
\[
       \operatorname{Tr}_{H_0}K_g
          =-\frac1{2\pi}\int_{\mathbb R}
                 G(it)\frac{m_\chi'(it)}{m_\chi(it)}\,dt
          =\mathcal W_{S,\chi}(g).
 \tag{MA-GA.25}
\]
This establishes the local trace from the actual operator
before any use of its scalar diagonal.  A constant phase
does not alter this trace; an actual exponential monomial
adds a constant multiple of $g(0)$, which is zero at the
asserted limit scope.  Neither is deleted from the operator.

\emph{5. The original finite-cutoff nuclear lift.}
For a fixed annulus $a=a_{\epsilon,X}$, the operator
$M_aC_g$ has the smooth compactly supported kernel
$a(x)g(x-y)$ on $\mathbb R^2$.  It maps $H_0$ into
$C_c^\infty(I)$ for one fixed compact output interval
$I$, continuously in every smooth seminorm.  It is
nuclear into that Fr\'echet space, not only trace class
into $H_0$.  To see the stronger assertion, choose an
interval containing the input support with an interior
margin, expand in its orthonormal Fourier basis in $y$,
and call the resulting compact-output coefficient
functions $v_j(x)$.  Repeated integration by parts in
$y$ gives
\[
       \|v_j\|_{C^k}\leq C_{k,N}(1+|j|)^{-N}
                         \quad\text{for every }k,N.
 \tag{MA-GA.26}
\]
There are positive numbers $d_j$ with $\sum_jd_j<\infty$
such that $\{v_j/d_j\}$ is bounded in every $C^k$ seminorm.
For example, choose successively increasing index blocks
on whose tails the first $r$ seminorms are at most
$2^{-r}(1+|j|)^{-r}$, and on the $r$th block use
$d_j=(1+|j|)^{-r/2}$, enlarging the finitely many initial
values; the block endpoints may be increased to make
the total sum finite.  Each fixed seminorm is then
bounded on all later blocks, as well as on the finitely
many earlier ones.

Thus $M_aC_g=\sum_jd_j(v_j/d_j)\otimes\ell_j$, where
the Fourier coefficient functionals $\ell_j$ have
uniformly bounded norm on $H_0$.  Composing them with
the bounded contraction $Q_{\chi,Y}$ gives one nuclear
series for $T^\chi_{X,Y,\epsilon}(g):H_0\to C_c^\infty(I)$.
The vector family is bounded in every output smooth
seminorm, and the functional bounds are uniform in $Y$
when the annulus and $g$ are fixed.

Compose its vector terms with (MA-GA.15).  They are now
bounded in one original compact, punctured LF stage of
$V_S$, flat at every coordinate hyperplane.  This is
the promised continuous nuclear return
\[
 \mathcal A_{X,Y,\epsilon}:H_\chi\longrightarrow V_S,
       \qquad \mathcal A_{X,Y,\epsilon}
                         =\sigma_\chi T^\chi_{X,Y,\epsilon}(g).
 \tag{MA-GA.27}
\]
The section alone is only asserted continuous; nuclearity
is asserted after this smoothing factor.  Let
$j_\chi:V_S\to H_\chi$ be centered periodization followed
by compact-character projection and the inclusion of
$\mathcal E$ in $H_\chi$.  Then
\[
 j_\chi\mathcal A_{X,Y,\epsilon}
            =T^\chi_{X,Y,\epsilon}(g),\qquad
 \mathcal A_{X,Y,\epsilon}j_\chi
          \text{ is a nuclear endomorphism of }V_S,
\]
\[
 \operatorname{Tr}_{V_S}(\mathcal A_{X,Y,\epsilon}j_\chi)
          =\operatorname{Tr}_{H_\chi}T^\chi_{X,Y,\epsilon}(g).
 \tag{MA-GA.28}
\]
Both traces use exactly the preceding summable series:
the scalar terms agree because $j_\chi\sigma_\chi=1$.
They are independent of its representation.  One can
verify the needed approximation property on the original
source directly: cut off each archimedean log coordinate
as in step 1, project onto a fixed finite-place stage,
and approximate on the resulting punctured compact
coordinate boxes by smooth finite Fourier sums with
a surrounding compact cutoff.  These finite-rank maps
converge to the identity uniformly on bounded sets in
each source seminorm; higher weights in (MA-GA.8)
control the first cutoff error.  This also supplies a
direct finite-rank justification of trace cyclicity.

\emph{6. Exact physical scaling and the uniform tails.}
Put $(W_RF)(x)=F(x+R)$.  Norm reversal by the actual
Fourier transform gives
\[
 W_RQ_{\chi,Y}W_{-R}=Q_{\chi,e^RY},\qquad
 W_{-\log Y}T^\chi_{X,Y,\epsilon}(g)W_{\log Y}
                       =T^\chi_{XY,1,\epsilon Y}(g).
 \tag{MA-GA.29}
\]
This is the same scaling of the two original physical
cutoffs as in MA-LA-014.  Use $U=XY$, $V=\epsilon Y$,
$a=A(e^x/U)-A(e^x/V)$, and (MA-GA.21).  The compact
kernel of $M_aC_gM_{b_1(-x)}$ has trace
$g(0)\int a(x)b_1(-x)\,dx$.  Therefore the exact trace
identity, including its subtraction at the identity, is
\[
 \operatorname{Tr}T^\chi_{X,Y,\epsilon}(g)
     =g(0)\int_{\mathbb R}a(x)b_1(-x)\,dx
                         +\operatorname{Tr}(M_aK_g).
 \tag{MA-GA.30}
\]
The first term vanishes at the stated $g(0)=0$ scope.
No assumption that the support of $g$ avoids zero is made.

Fix $0<\delta<1/2$.  The multiplication map
$1-A(e^x/U):H_\delta\to H_0$ has norm at most
$U^{-\delta}$, since its support is in $x\geq\log U$.
The map $A(e^x/V):H_{-\delta}\to H_0$ has norm at most
$(2V)^\delta$, since its support is in
$x\leq\log(2V)$.  Equations (MA-GA.24--25) and
$1-a=1-A(e^x/U)+A(e^x/V)$ yield
\[
 \bigl|\operatorname{Tr}(M_aK_g)-\operatorname{Tr}K_g\bigr|
       \leq U^{-\delta}\|K_g:H_0\to H_\delta\|_1
            +(2V)^\delta\|K_g:H_0\to H_{-\delta}\|_1.
 \tag{MA-GA.31}
\]
Finally $U\geq2V$ implies $0\leq a\leq1$ and
$\operatorname{supp}a\subset[\log V,\log(2U)]$.
Consequently
\[
      |\operatorname{Tr}(M_aK_g)|
           \leq V^{-\delta}\|K_g:H_0\to H_\delta\|_1.
 \tag{MA-GA.32}
\]
These inequalities prove all the asserted rates and limits.
They also specify exactly the uniformity that is required;
no limiting interchange over compact characters, unit
lattices, fields or finite sets of places has been used.
\end{proof}

\paragraph{Source locators and interpretation.}
Masullo's public draft, privately downloaded for this run,
defines the flat-hyperplane
space in Lemma 5.3.3.3(3), printed/PDF pp.~252--257,
especially (5.3.3.10), and the $S$-local source in
Definition 5.3.4.1, p.~259, (5.3.4.1), with its LF
properties discussed in Lemma 5.3.4.2, pp.~260--261.
The direct Taylor proof (MA-GA.8) is the proof used here
for the needed topology and density estimates.

Meyer (2005), \S\S3.1--3.2, supplies the $S$-unit quotient
and the actual unitary descent; exact supplied-source
locators are \texttt{eq:def\_Fourier}, \texttt{eq:def\_J},
\texttt{eq:def\_EFou}, and
\texttt{lem:Fourier\_unitary\_ICL}, including its full proof.
The weighted extension is Proposition 4.4,
\texttt{pro:Fourier\_estimate}.  The original full source
identification is the complete argument ending at
\texttt{cor:Sch\_local\_coinvariants}, preceded by
\texttt{lem:S\_local\_temperation} and
\texttt{the:Sch\_Adel\_temperation}; the structural condition
is the sufficiently-large-$S$ definition preceding the
global construction.  These locators do not identify
the flat quotient with a different bornology: that
additional point is proved in (MA-GA.8--15).

The analytic local forms are Tate (1950), Theorem 2.4.1,
thesis p.~(2.11), and \S2.5, real p.~(2.14), complex
p.~(2.16), nonarchimedean p.~(2.19), respectively PDF
pages 18, 21, 23 and 26 of the original-thesis scan.
Formula (MA-GA.17) fixes their use to the actual chosen
Fourier convention.  The scalar trace is also Meyer's
general-character local trace formula (20), p.~22;
it was derived here from the actual kernel before being
identified with that local distribution.  In the rational
even sector it gives $A_0(h)+B_S(h)$ of MA-LA-014 for
$g(x)=h(e^x)$; in the sectors of MA-CQ it gives precisely
their ramified and archimedean local terms.  The fixed
compact-character projector has trace one, so no extra
class-number or regulator multiplicity is inserted.

The nuclear comparison is an actual source-to-quotient
comparison with the common pairing (MA-GA.28).  It
does not assert the separate global-to-local identity
FT002, or any positivity or location of its discarded
zeros.  For finite-order characters it retains the same
finite-monodromy sector as MA-CQ; allowing a general
unitary compact character does not assign it an
unconstructed finite Artin motive.
